\documentclass{article}
\usepackage[utf8]{inputenc}
\usepackage{a4wide}

\begin{document}

\section*{ธุรกรรมอำพราง}

ในการจัดการบัญชีรวมของลูกค้าคนหนึ่งของธนาคารจะมีทั้งบัญชีเงินฝาก และบัญชีเงินกู้

โดยลูกค้าทุกคนสามารถทำธุรกรรมด้านการฝากเงิน ถอนเงิน หรือกู้เงินจากบัญชีรวมนี้ได้

โดยเป้าหมายของธนาคาร

คือไม่ต้องการให้ลูกค้าคนใดคนหนึ่งมียอดเงินคงเหลือในบัญชีรวมธุรกรรมนี้เป็น 0

ในช่วงใดช่วงหนึ่งของการทำธุรกรรม เพราะนั้นหมายถึงลูกค้าคนนั้นไม่ทั้งยอดเงินฝาก

และยอดเงินกู้ ส่งผลให้ลูกค้าอาจจะปิดบัญชี และย้ายไปเป็นลูกค้าของธนาคารอื่นได้



ลูกค้าคนหนึ่งได้ทำธุรกรรมทั้งหมด n ครั้ง

ในแต่ละครั้งของการทำธุรกรรมจะแสดงยอดสรุปของบัญชีรวม ณ ขณะนั้น

ถ้ายอดเป็นเลขบวกหมายถึงมีเงินในบัญชีเงินฝากมากกว่ายอดในบัญชีเงินกู้

ในทางกลับกันถ้ายอดสรุปเป็นลบหมายถึงยอดในบัญชีเงินกู้มากกว่าบัญชีเงินฝาก

(ยอดสรุปนี้จะไม่สามารถเป็น 0 ได้เพราะธนาคารจะไม่ยอมให้เกิดขึ้น)

จากการวิเคราะห์ข้อมูลทางธนาคารพบว่าธนาคารมักจะเสียลูกค้าที่มีการทำกลุ่มของธุรกรรมย่อยที่ติดกัน

และส่งผลให้ผลรวมของยอดสรุปของกลุ่มธุรกรรมย่อยที่ติดกันนั้นมีค่าเป็น 0

ซึ่งถือเป็นปัญหาสำหรับธนาคาร โดยธนาคารตัดสินใจเพิ่มธุรกรรมเสมือนซึ่งเป็นค่าตัวเลขอะไรก็ได้

เพื่อไม่ให้ผลรวมของยอดสรุปของกลุ่มธุรกรรมย่อยที่ติดกันนั้นมีค่าเป็น 0

(ไม่มีผลต่อยอดเงินรวมทั้งหมดของลูกค้า)



\ul{งานของคุณ}:

ช่วยธนาคารหาจำนวนครั้งที่น้อยที่สุดที่ต้องแทรกการทำธุรกรรมเสมือนระหว่างการทำธุรกรรมเดิมของลูกค้า

เพื่อให้มั่นใจว่าลูกค้าจะไม่มีผลรวมยอดเงินของธุรกรรมย่อยที่ติดกันใดๆ เป็น 0



\hfill\break



\hfill\break



\hfill\break



{\def\LTcaptype{none} % do not increment counter

\begin{longtable}[]{@{}

  >{\raggedright\arraybackslash}p{(\linewidth - 4\tabcolsep) * \real{0.3333}}

  >{\raggedright\arraybackslash}p{(\linewidth - 4\tabcolsep) * \real{0.3333}}

  >{\raggedright\arraybackslash}p{(\linewidth - 4\tabcolsep) * \real{0.3333}}@{}}

\toprule\noalign{}

\begin{minipage}[b]{\linewidth}\centering

ข้อมูลเข้า

\end{minipage} & \begin{minipage}[b]{\linewidth}\centering

ข้อมูลออก

\end{minipage} & \begin{minipage}[b]{\linewidth}\centering

คำอธิบาย

\end{minipage} \\

\midrule\noalign{}

\endhead

\bottomrule\noalign{}

\endlastfoot

\begin{minipage}[t]{\linewidth}\raggedright

5\\

3 -6 4 2\\

\strut

\end{minipage} & 1 & \begin{minipage}[t]{\linewidth}\raggedright

ในตัวอย่างนี้ มีช่วงย่อยที่ผลรวมเป็น 0 เพียงช่วงเดียว โดยเริ่มจากตัวที่ 2 ถึงตัวที่ 4\\

ดังนั้นจึงจำเป็นต้องแทรกจำนวนหนึ่งระหว่างตัวที่ 2 และ 3 เพื่อให้ไม่มีช่วงย่อยที่ผลรวมเป็น 0\\

\strut

\end{minipage} \\

\begin{minipage}[t]{\linewidth}\raggedright

4\\

3 -2 1 -8\\

\strut

\end{minipage} & 0 & \begin{minipage}[t]{\linewidth}\raggedright

ในกรณีนี้ไม่มีช่วงย่อยที่ผลรวมเป็น 0 ดังนั้นไม่จำเป็นต้องทำอะไร\\

\strut

\end{minipage} \\

\begin{minipage}[t]{\linewidth}\raggedright

7\\

-1 1 -1 1 1 -1 -1\strut

\end{minipage} & 4 & \begin{minipage}[t]{\linewidth}\raggedright

ในกรณีนี้มีช่วงย่อยหลายช่วงที่ผลรวมเป็น 0

และจำเป็นต้องแทรกจำนวนเพื่อแยกช่วงย่อยเหล่านี้ออกจากกัน โดยต้องแทรกจำนวนทั้งหมด 4

ตัว\\

\strut

\end{minipage} \\

\begin{minipage}[t]{\linewidth}\raggedright

6\\

12 -5 -7 22 4 -4\strut

\end{minipage} & 2 & \begin{minipage}[t]{\linewidth}\raggedright

ในกรณีนี้มีช่วงย่อยที่ผลรวมเป็น 0 อยู่ 2 ช่วง ดังนั้นจำเป็นต้องแทรกจำนวน 2

ตัวเพื่อไม่ให้มีช่วงย่อยที่ผลรวมเป็น 0\\

\strut

\end{minipage} \\

\end{longtable}

}



\subsection*{Input}
บรรทัดแรกประกอบด้วยจำนวนเต็ม n (2 ≤ n ≤ 200000) --- จำนวนธุรกรรมของลูกค้า



บรรทัดที่สองประกอบด้วยจำนวนเต็ม n ตัว a1, a2, ..., an (−10\^{}9 ≤ ai ≤

10\^{}9, ai ≠ 0) --- ยอดเงินของแต่ละธุรกรรม\\



\subsection*{Output}
จำนวนเต็มหนึ่งตัว --- จำนวนครั้งที่น้อยที่สุดที่ธนาคารต้องสร้างธุรกรรมเสมือน

เพื่อให้ไม่มีผลรวมยอดเงินของธุรกรรมย่อยที่ติดกันใดๆ เป็น 0\\



\end{document}
